\phantomsection
\chapter*{Introducción}
\addcontentsline{toc}{chapter}{Introducción}

La educación en la era digital demanda habilidades que trascienden la lectura, la escritura y las matemáticas tradicionales. Entre estas, el pensamiento computacional se ha consolidado como una aptitud fundamental para la resolución de problemas, el diseño de sistemas y la comprensión del mundo tecnológico que rodea a las nuevas generaciones. Su desarrollo desde edades tempranas resulta esencial para preparar a niños y jóvenes como creadores y pensadores críticos capaces de desenvolverse en una sociedad donde los conceptos de computación están siempre presentes.

Sin embargo, las herramientas educativas existentes para fomentar el pensamiento computacional, como Scratch, se fundamentan exclusivamente en el paradigma de flujo de control, lo que favorece la formación de modelos mentales rígidamente secuenciales. Esta limitación dificulta la posterior comprensión de conceptos avanzados como la concurrencia, cada vez más relevante en la computación moderna. La ausencia de alternativas basadas en otros paradigmas de programación restringe las posibilidades de formación y reduce la diversidad de medios para el desarrollo de esta habilidad crítica en el contexto educativo actual.

Ante esta problemática, el presente trabajo de grado tiene como propósito el desarrollo de NodeLand, un lenguaje de programación visual basado en el paradigma de flujos de datos (dataflow), y su entorno, diseñado específicamente para hacer accesible, intuitiva y atractiva la formación del pensamiento computacional en niños y jóvenes de 8 a 16 años. La propuesta se sustenta teóricamente en las definiciones de pensamiento computacional de Wing y Aho, el marco de conceptos, prácticas y perspectivas de Brennan y Resnick, las etapas de desarrollo cognitivo de Piaget, y el construccionismo de Papert. Metodológicamente, el proyecto se enmarca como una investigación proyectiva, y su desarrollo se llevó a cabo siguiendo un modelo incremental de software en cuatro iteraciones que abarcaron desde la construcción del editor visual y el motor de ejecución concurrente, pasando por la integración con hardware micro:bit y computadoras remotas, hasta la incorporación de un módulo de desafíos gamificados y una plataforma social colaborativa, validándose mediante evaluación experta y pruebas de usabilidad con el público objetivo.

Este trabajo se encuentra organizado en cinco capítulos. En el Capítulo I se aborda la problemática que dio origen al estudio, junto con la formulación de los objetivos, el alcance definido, las limitaciones identificadas y la justificación que respalda la investigación. El Capítulo II desarrolla los antecedentes y el marco teórico, en el cual se incluyen las bases conceptuales sobre pensamiento computacional, fundamentos cognitivos y pedagógicos, lenguajes de programación, programación por flujos de datos, concurrencia, sistemas distribuidos y micro:bit. El Capítulo III describe el enfoque metodológico adoptado, precisando el tipo de investigación, las técnicas e instrumentos de recolección de datos y la metodología de desarrollo utilizada. En el Capítulo IV se expone el desarrollo y los resultados obtenidos, detallando el análisis de la problemática, el diseño del lenguaje y su entorno, la implementación incremental, la validación con expertos y usuarios del público objetivo, y la documentación elaborada. Por último, el Capítulo V presenta las conclusiones derivadas del trabajo y plantea recomendaciones para futuras líneas de investigación.